% =====================================================================
%  Método 6 — Teste do núcleo predicativo   (dimensão: estrutura)
%  Origem: células 25, 26 e 27 do caderno
% =====================================================================
\subsection{Teste do núcleo predicativo}
\label{met:raiz}

Toda frase requer um predicado. Não se trata de heurística nem de tendência estatística, mas
de propriedade sintática, e em uma árvore de dependências essa propriedade é verificável em
um único ponto: a raiz, que por construção é o nó do qual todos os demais dependem.

Seja $\raiz(\sent)$ a raiz da árvore e $\mathrm{filhos}(\raiz)$ seus dependentes diretos. A
frase possui predicado quando
%
\begin{equation}
\mathrm{Pred}(\sent) \;=\;
\underbrace{\bigl[\mathrm{pos}(\raiz) \in \{\texttt{VERB}, \texttt{AUX}\}\bigr]}_{\text{predicado verbal}}
\;\lor\;
\underbrace{\bigl[\exists\, c \in \mathrm{filhos}(\raiz) : \mathrm{dep}(c) = \texttt{cop}\bigr]}_{\text{predicado nominal}}
\end{equation}
%
Duas linhas de código e um único nó inspecionado bastam, o que faz deste o método de menor
custo entre os onze e ainda assim produz veredito binário e localizado.

A inspeção recai sobre a raiz, e não sobre a presença de verbo, porque o teste ingênuo ---
verificar se existe algum verbo na frase --- aprova o fragmento. Confrontem-se \emph{O
incêndio \textbf{atingiu} a mata nativa} e \emph{O incêndio que \textbf{atingiu} a mata
nativa}: ambas contêm o verbo \texttt{atingiu}, mas a primeira constitui oração e a segunda
não. O que difere não é a presença do verbo, e sim sua posição na árvore, pois o pronome
\texttt{que} o rebaixa a oração relativa (\texttt{acl:relcl}) e promove o substantivo à raiz.
Uma única palavra converte oração em fragmento, e apenas a estrutura registra o fato.

O segundo termo da disjunção é igualmente necessário. Em predicações nominais, como \emph{O
carro é azul}, a raiz é o predicativo \texttt{azul}, etiquetado \texttt{ADJ}, e o verbo
\texttt{é} figura como dependente com função \texttt{cop}. Na ausência dessa cláusula, toda
construção copular seria classificada como fragmento, o que inviabilizaria o método em
português.

O método devolve um veredito binário acompanhado do nó que o justifica, conforme a
Tabela~\ref{tab:raiz}.

\begin{table}[H]
\centering
\dimtable{
\begin{tabular}{@{}llll@{}}
\toprule
\textbf{Frase} & \textbf{Raiz} & \textbf{Classe} & \textbf{Veredito} \\
\midrule
Há três dias o incêndio atingiu uma área\ldots       & \texttt{atingiu} & VERB & ORAÇÃO \\
Há três dias o incêndio que atingiu uma área\ldots   & \texttt{dias}    & NOUN & \textbf{FRAGMENTO} \\
O incêndio é intenso.                                & \texttt{intenso} & ADJ + cópula & ORAÇÃO \\
\bottomrule
\end{tabular}}
\caption{O teste da raiz sobre três frases do corpus.}
\label{tab:raiz}
\end{table}

A segunda linha corresponde ao fragmento, e a terceira exercita a cláusula da cópula: a raiz é
um adjetivo, o primeiro termo da disjunção é falso, e ainda assim a frase é corretamente
aprovada. O método mensura completude, e nada além disso, não estabelecendo nada quanto à
correção ou à plausibilidade da frase.
